\section{Related Work}
\label{sec:related}\label{sec:relatedwork}
\paragraph{Quantisation and grouped ADC.}
PQ~\cite{pq} decomposes vectors into subspaces for compressed storage and
asymmetric distance computation. JHQ~\cite{jhq} supplies our transform,
Cartesian primary codebook and residual hierarchy. PQ Fast
Scan~\cite{fastscan} and Quicker ADC~\cite{quickeradc} establish small-table
execution, including 16-entry tables. We study grouping and packed access for
JHQ's representation. 
%Equation~\ref{eq:split} follows from its Cartesian
%construction; arbitrary learned PQ subcodebooks lack this general guarantee
%of internal separability.

\paragraph{GPU quantised search.}
IVF-RaBitQ~\cite{ivfrabitq} is the closest execution design: a cheap filter
feeding a refinement stage, grouped lookup tables and a GPU layout. Two-stage
execution and grouped lookup are shared ideas; ours are their measured
trade-off for JHQ's Cartesian codebook and the calibration of its residual
budget. We do not report it as a baseline: the available implementation exposes
search and refinement separately, and we did not validate an assembled
two-stage pipeline against the published system. CAGRA~\cite{cagra} and cuVS
IVF-PQ~\cite{cuvs} supply the GPU comparisons; GPU compressed-domain scanning
and $k$-selection go back to Faiss-GPU~\cite{faissgpu}.

\paragraph{Adaptive search.}
Faiss~\cite{faisslib} sweeps search-time settings and keeps the
Pareto-optimal ones; learned early termination~\cite{earlyterm} controls
per-query search effort. Our workload
policy uses unlabelled output agreement to select JHQ's residual budget, with
enumeration as an empirical reference. It calibrates a shared budget for
subsequent batches under stationary queries; it neither predicts per-query
difficulty nor guarantees a target ground-truth recall.
